\chapter*{VI/20 Legacy Problem Recovery Audit}
\addcontentsline{toc}{chapter}{VI/20 Legacy Problem Recovery Audit}

\section*{Audit invariant}
The Volume VI invariant remains
\[
\boxed{\text{every substantive legacy problem}\longrightarrow\text{exactly one canonical destination}.}
\]
VI/20 requires a correction to the pre-existing ledger: the direct AG3 gluing problem was already reserved here, but AG5 Problem 2 was routed here by the scope audit without its own legacy-ledger row.

\section*{Corpus centers assigned to VI/20}
Two substantive numbered legacy problems belong to this chapter.
\begin{center}
\begin{tabular}{>{\raggedright\arraybackslash}p{0.14\textwidth}>{\raggedright\arraybackslash}p{0.43\textwidth}>{\raggedright\arraybackslash}p{0.28\textwidth}}
\toprule
Corpus item & Mathematical center & VI/20 decision\\
\midrule
AG3 Problem 14 & General gluing theorem for schemes along compatible open-subscheme isomorphisms; quotient topology, glued sheaf, local stalks, scheme verification, uniqueness & INCLUDE as VI.20.P1; preserved in full\\
AG5 Problem 2 & Intersections of affine open subschemes are covered by affine opens distinguished in both charts & INCLUDE as VI.20.P2; restore omitted ledger entry\\
AG3 Problem 13 & $\Spec\mathbb Z$ final / structure morphism & KEEP RESERVED for VI/21\\
\bottomrule
\end{tabular}
\end{center}

\section*{AG3 Problem 14 preservation check}
VI.20.P1 retains every substantive step in the source solution:
\begin{itemize}
\item extension by $U_{ii}=X_i$ and identity transitions;
\item equivalence relation on the disjoint union, with inverse and cocycle proofs;
\item quotient/final topology and openness of each chart image;
\item injectivity/homeomorphism of charts onto their images;
\item exact description of chart intersections and transition compatibility;
\item definition of $\OO_X(V)$ as compatible families of chart sections;
\item componentwise restriction and the sheaf proof;
\item recovery of each original chart sheaf on its image;
\item identification of glued stalks with local chart stalks;
\item verification that the result is a scheme via the open chart cover;
\item uniqueness up to unique chart-compatible isomorphism.
\end{itemize}
The canonical mapping-property, two-chart, equalizer, doubled-origin, and reconstruction problems are derivatives/applications of this one legacy theorem, not further migrations of AG3-P14.

\section*{AG5 Problem 2 recovery and ledger repair}
The scope audit VI-SCOPE-038 assigns the affine-intersection/gluing descendants in the AG5 source to VI/20.  The source contains a substantive numbered
``Problem 2 — Intersections of Affine Open Subschemes.''  The old legacy ledger jumped directly to
AG5-P3 and therefore silently omitted this item.  VI.20.P2 now preserves the full shrinking argument:
principal basis in the first affine chart, refinement in the second, localization inside a principal
open, and a final common distinguished refinement.  A new AG5-P2 ledger row is inserted rather than
hiding the problem inside theory prose.

\section*{Nearby variants and non-duplication}
The chapter develops useful canonical consequences of P1/P2: two-chart gluing, the gluing mapping
property, reconstruction from affine covers, principal-open ring gluing, equalizers of global sections,
the doubled-origin line, inversion gluing, associativity, atlas refinement, and compatible morphism
gluing.  These are explicitly canonical derivatives.  They do not create second legacy destinations.

VI/18's $D(f)\cong\Spec A_f$ theorem and VI/19's affine anti-equivalence are used only as prior
infrastructure to express overlap schemes algebraically; neither is re-migrated.

\section*{Boundary audit}
\begin{itemize}
\item AG3-P13 remains \texttt{RESERVED} for VI/21; VI/20 does not prove finality of $\Spec\mathbb Z$ as a canonical dossier.
\item Open/closed subscheme formalism remains centered in VI/22.  VI/20 uses open subschemes only as the interface already required by AG3-P14.
\item Fiber products and base change remain reserved for VI/23--VI/24.
\item The two-chart inversion example is not counted as a migration of later projective-space material; its projective identification is explicitly deferred.
\end{itemize}

\section*{No-loss status after VI/20}
The ledger/scope state after this package is:
\begin{itemize}
\item AG3-P14 $\to$ VI.20.P1, status \texttt{MIGRATED\_DRAFT};
\item AG5-P2 $\to$ VI.20.P2, status \texttt{MIGRATED\_DRAFT}, newly inserted into the legacy ledger;
\item VI-SCOPE-038 receives \texttt{VI20\_AFFINE\_GLUING\_INSTANCE\_AUDIT\_2026-08-21};
\item AG3-P13 remains \texttt{RESERVED} for VI/21.
\end{itemize}
Thus the gluing corpus is accounted for without dropping the affine-intersection problem or duplicating the general AG3 gluing theorem.
